%==============================================================================
%  MATH 231 -- Linear Algebra I (Queens College, Fall 2026)
%  MATRIX UNIT, PART 11 -- Gram-Schmidt and the QR Decomposition
%
%  Section numbering runs CONTINUOUSLY across the parts of the matrix unit;
%  see Part 2's header for the full map.  This part is Sections 55-58, hence
%  \setcounter{section}{54}.
%
%  STATUS: complete.  Sections 55-58.  Written from the plan in
%  plan/MATH231_remaining_units_plan.tex (2026-10-25).  This is the last
%  part of the matrix unit, and the last topic for Midterm 2.
%
%  AUDIENCE: distributed to students.  No instructor-facing apparatus.
%
%  WHAT THE INSTRUCTOR SPECIFIED, AND WHERE IT IS.
%   - The Gram-Schmidt algorithm/process: Sec. 55.
%   - The QR decomposition: Sec. 56.
%   - col[Q] = col[A]: Sec. 57.
%   - Applying QR to least squares: Sec. 58.
%
%  DECISIONS CARRIED FROM THE PLAN.
%
%   D9. Gram-Schmidt is derived in full in computational unit Part 4
%       (Sec. 16-19), which may or may not have been lectured.  Sec. 55 is
%       written to stand on its own (the process, the pseudocode, why w_k is
%       never zero, the worked example) but compactly, citing the
%       computational unit for the derivation from the R^2 picture and for
%       the dependent case.  If computational Part 4 is lectured first,
%       Sec. 55 can be read as a recap.
%
%   DEBT COLLECTED.  Computational unit Sec. 20.3 closes: "the relationship
%   between the input list and the output list ... is itself a factorisation
%   of the matrix whose columns are the v_k.  Naming it requires matrix
%   machinery we do not yet have.  We will return to Gram-Schmidt."
%   Sec. 56.5 says this is that return.
%
%   NOTATION.  The computational unit writes the inputs v_1, ..., v_n; here
%   they are the columns a_1, ..., a_n of A, and Sec. 55.1 says so.  The
%   pseudocode is CLASSICAL Gram-Schmidt, as in the computational unit, and
%   Sec. 58.4 points at the modified version (computational Sec. 20.3) and
%   at Householder reflections, which are built in Sec. 28.5-28.7.
%
%   WORKED EXAMPLES reuse computational Sec. 18.3's vectors [1,1,0],
%   [1,0,1], [0,1,1] (q's already known there) and the line-fit data of
%   Sec. 49.4.
%
%  Compile with pdflatex or tectonic (standard TeX Live packages only).
%==============================================================================
\documentclass[11pt]{article}

\usepackage[margin=1in]{geometry}
\usepackage{amsmath,amssymb,amsthm}
\usepackage[dvipsnames]{xcolor}
\usepackage{enumitem}
\usepackage{booktabs}
\usepackage{array}
\usepackage[hidelinks]{hyperref}
\usepackage{titlesec}
\usepackage{needspace}

%------------------------------------------------------------------ style ----
\titleformat{\section}{\large\bfseries\sffamily}{\thesection.}{0.6em}{}
\titleformat{\subsection}{\normalsize\bfseries\sffamily}{\thesubsection.}{0.5em}{}
\titlespacing*{\section}{0pt}{14pt}{6pt}
\setlength{\emergencystretch}{3em}

\theoremstyle{plain}
\newtheorem{thm}{Theorem}[section]
\newtheorem{lem}[thm]{Lemma}
\newtheorem{cor}[thm]{Corollary}

\theoremstyle{definition}
\newtheorem{defn}{Definition}[section]
\newtheorem{ex}{Example}[section]
\theoremstyle{remark}
\newtheorem*{rmk}{Remark}
\newtheorem*{warn}{A common confusion}

\newcommand{\lookahead}{\par\smallskip\noindent
  \textbf{\sffamily\textcolor{RoyalBlue}{Looking ahead.}}\ \ignorespaces}
\newcommand{\lookback}[1][the vector unit]{\par\smallskip\noindent
  \textbf{\sffamily\textcolor{ForestGreen}{From #1.}}\ \ignorespaces}

%--------------------------------------------------------------- notation ----
\newcommand{\R}{\mathbb{R}}
\newcommand{\vc}[1]{\mathbf{#1}}
\newcommand{\mt}[1]{\mathbf{#1}}
\newcommand{\norm}[1]{\lVert #1\rVert}
\newcommand{\T}{^{\mathsf{T}}}
\newcommand{\colsp}[1]{\operatorname{col}[#1]}
\newcommand{\xh}{\hat{\vc{x}}}
\DeclareMathOperator{\spanop}{span}
\DeclareMathOperator{\rank}{rank}

%--------------------------------------------------------------- pseudocode --
\newcommand{\kw}[1]{\textsf{\bfseries #1}}
\newcommand{\ind}{\hspace*{1.3em}}
\newcommand{\indd}{\hspace*{2.6em}}
\newcounter{algoline}
\newenvironment{algo}[2]{%
  \par\medskip\needspace{8\baselineskip}%
  \begin{center}\begin{minipage}{0.94\textwidth}%
  \hrule height 0.9pt \vspace{3pt}%
  \noindent{\sffamily\bfseries #1}\par\vspace{2pt}%
  \noindent{\small #2}\par\vspace{3pt}%
  \hrule height 0.4pt \vspace{5pt}%
  \begin{list}{\textbf{\arabic{algoline}:}}%
    {\usecounter{algoline}\setlength{\leftmargin}{2.3em}%
     \setlength{\labelwidth}{1.8em}\setlength{\labelsep}{0.5em}%
     \setlength{\itemsep}{2pt}\setlength{\topsep}{0pt}%
     \setlength{\parsep}{0pt}}%
}{%
  \end{list}\vspace{5pt}\hrule height 0.9pt%
  \end{minipage}\end{center}\medskip%
}

\title{\vspace{-1.2cm}\textbf{\sffamily MATH 231 -- Linear Algebra I}\\[4pt]
       {\large\sffamily The Matrix Unit, Part 11 \textbullet\ Gram--Schmidt
        and the $QR$ Decomposition}\\[2pt]
       {\normalsize\sffamily Kiely Hall 326 \textbullet\
        Instructor: Kevin Bedoya}}
\author{}
\date{}

\begin{document}
\setcounter{section}{54}  % this part opens at Section 55; see the header note
\maketitle
\vspace{-1.6cm}

\noindent\textbf{About this document.} This is Part 11, the last part of the
matrix unit: Sections 55--58. Section numbers run continuously through the unit
(Part 2's opening table lists them all). References into the computational unit
always name it.

\medskip
\noindent\textbf{Assumed.} Orthonormal vectors and orthogonal matrices (\S20);
projections $\mt{P}_{\vc{u}}=\vc{u}\vc{u}\T/\vc{u}\T\vc{u}$ (\S26); the column
space (\S10); triangular matrices, their determinants and inverses (\S39, \S53);
back substitution (\S41); least squares and the pseudo-inverse (Part 9, \S54).
The Gram--Schmidt process is derived in the computational unit, Part 4
(\S16--\S19); \S55 below restates what is needed.

\medskip
\noindent\textbf{What you should be able to do afterwards.} Run the Gram--Schmidt
process on a few vectors; rewrite its output as a factorization
$\mt{A}=\mt{Q}\mt{R}$ and compute $\mt{R}$; prove that $\colsp{\mt{Q}}=\colsp{\mt{A}}$;
solve a least-squares problem with $\mt{Q}$ and $\mt{R}$ by one back
substitution; and say why this is preferred to the normal equations in practice.

%==============================================================================
\section{Gram--Schmidt, in matrix language}
%==============================================================================

%------------------------------------------------------------------------------
\subsection{The goal}
%------------------------------------------------------------------------------
Let $\mt{A}\in\R^{m\times n}$ have independent columns $\vc{a}_1,\dots,\vc{a}_n$
(so $m\ge n$). The \textbf{Gram--Schmidt process} produces orthonormal vectors
$\vc{q}_1,\dots,\vc{q}_n$ that span the same spaces, step by step:
\[
  \spanop\{\vc{q}_1,\dots,\vc{q}_k\}=\spanop\{\vc{a}_1,\dots,\vc{a}_k\}
  \qquad\text{for every }k=1,\dots,n .
\]
\lookback[the computational unit] Part 4 built the process from a single move in
$\R^{2}$, ``subtract the projection, keep what is orthogonal'', and named its
inputs $\vc{v}_k$. Here they are the columns $\vc{a}_k$ of a matrix, because the
point of this part is what the process says about that matrix.

%------------------------------------------------------------------------------
\subsection{The process}
%------------------------------------------------------------------------------
Take the columns in order. From each, subtract its projections onto the
$\vc{q}$'s already built, then normalize what remains:
\[
  \vc{w}_k=\vc{a}_k-\sum_{j=1}^{k-1}(\vc{q}_j\T\vc{a}_k)\,\vc{q}_j,
  \qquad
  \vc{q}_k=\frac{\vc{w}_k}{\norm{\vc{w}_k}} .
\]
Each subtracted term is a projection: since $\vc{q}_j$ is a unit vector,
$(\vc{q}_j\T\vc{a}_k)\vc{q}_j=\vc{q}_j\vc{q}_j\T\vc{a}_k=\mt{P}_{\vc{q}_j}\vc{a}_k$
(\S26.2). Because the earlier $\vc{q}_j$ are orthonormal, $\vc{w}_k$ is orthogonal
to each of them: for $i<k$,
$\vc{q}_i\T\vc{w}_k=\vc{q}_i\T\vc{a}_k-\sum_j(\vc{q}_j\T\vc{a}_k)\,\vc{q}_i\T\vc{q}_j=\vc{q}_i\T\vc{a}_k-\vc{q}_i\T\vc{a}_k=0$.

\textbf{$\vc{w}_k$ is never zero.} By construction $\vc{a}_k=\vc{w}_k+(\text{a
combination of }\vc{q}_1,\dots,\vc{q}_{k-1})$, and those $\vc{q}$'s span the same
space as $\vc{a}_1,\dots,\vc{a}_{k-1}$. If $\vc{w}_k$ were $\vc{0}$, $\vc{a}_k$
would be a combination of the earlier columns, contradicting independence. So
the division is always safe. What happens when the columns \emph{are} dependent
is the computational unit's \S19.

\begin{algo}{Gram--Schmidt (classical)}{%
  \textbf{Input:} independent columns $\vc{a}_1,\dots,\vc{a}_n\in\R^{m}$.\quad
  \textbf{Output:} orthonormal $\vc{q}_1,\dots,\vc{q}_n$, and the numbers
  $r_{jk}$ of \S56.}
\item \kw{for} $k = 1,\dots,n$:
\item \ind $\vc{w} \;\gets\; \vc{a}_k$
\item \ind \kw{for} $j = 1,\dots,k-1$:
\item \indd $r_{jk} \;\gets\; \vc{q}_j\T\vc{a}_k$
        \hfill\textit{\small the length of the shadow on $\vc{q}_j$}
\item \indd $\vc{w} \;\gets\; \vc{w}-r_{jk}\,\vc{q}_j$
        \hfill\textit{\small remove it}
\item \ind $r_{kk} \;\gets\; \norm{\vc{w}}$
\item \ind $\vc{q}_k \;\gets\; \vc{w}\,/\,r_{kk}$
\item \kw{return} $\vc{q}_1,\dots,\vc{q}_n$ and all $r_{jk}$
\end{algo}

\noindent The only change from the computational unit's pseudocode is that the
numbers the process computes along the way, the shadow lengths $r_{jk}$ and the
lengths $r_{kk}$, are now \emph{kept}. They are the second factor of the
factorization in \S56.

%------------------------------------------------------------------------------
\subsection{A worked example}
%------------------------------------------------------------------------------
\begin{ex}
\label{ex:gs}
The computational unit's \S18.3 example, as the columns of a matrix:
\[
  \mt{A}=\begin{bmatrix}1&1&0\\1&0&1\\0&1&1\end{bmatrix},
  \qquad
  \vc{a}_1=[\,1,1,0\,],\ \vc{a}_2=[\,1,0,1\,],\ \vc{a}_3=[\,0,1,1\,].
\]
\emph{$k=1$.} $r_{11}=\norm{\vc{a}_1}=\sqrt2$ and $\vc{q}_1=\frac{1}{\sqrt2}[\,1,1,0\,]$.

\emph{$k=2$.} $r_{12}=\vc{q}_1\T\vc{a}_2=\frac{1}{\sqrt2}$, and
$\vc{w}_2=[\,1,0,1\,]-\frac{1}{\sqrt2}\cdot\frac{1}{\sqrt2}[\,1,1,0\,]=[\,\frac12,-\frac12,1\,]$,
with $r_{22}=\norm{\vc{w}_2}=\sqrt{\frac32}=\frac{\sqrt6}{2}$. So
$\vc{q}_2=\frac{1}{\sqrt6}[\,1,-1,2\,]$.

\emph{$k=3$.} $r_{13}=\vc{q}_1\T\vc{a}_3=\frac{1}{\sqrt2}$ and
$r_{23}=\vc{q}_2\T\vc{a}_3=\frac{-1+2}{\sqrt6}=\frac{1}{\sqrt6}$, so
\[
  \vc{w}_3=[\,0,1,1\,]-\tfrac12[\,1,1,0\,]-\tfrac16[\,1,-1,2\,]
  =\bigl[\,-\tfrac23,\ \tfrac23,\ \tfrac23\,\bigr],
\]
with $r_{33}=\norm{\vc{w}_3}=\frac{2}{\sqrt3}$ and $\vc{q}_3=\frac{1}{\sqrt3}[\,-1,1,1\,]$.
These are the vectors the computational unit found.
\end{ex}

%==============================================================================
\section{The \texorpdfstring{$QR$}{QR} decomposition}
%==============================================================================

%------------------------------------------------------------------------------
\subsection{Running the process backwards}
%------------------------------------------------------------------------------
The process says how to get each $\vc{q}_k$ from the $\vc{a}$'s. Solve it the other
way, for $\vc{a}_k$. From $\vc{w}_k=r_{kk}\vc{q}_k$ and the definition of
$\vc{w}_k$,
\[
  \vc{a}_k=r_{1k}\vc{q}_1+r_{2k}\vc{q}_2+\cdots+r_{k-1,k}\vc{q}_{k-1}+r_{kk}\vc{q}_k
  =\sum_{j=1}^{k}r_{jk}\,\vc{q}_j .
\]
Column $k$ of $\mt{A}$ is a combination of the first $k$ $\vc{q}$'s only. By the
column reading (\S2.5), a combination of $\vc{q}_1,\dots,\vc{q}_n$ with weights
$r_{1k},\dots,r_{nk}$ is the matrix $\mt{Q}=[\,\vc{q}_1\ \cdots\ \vc{q}_n\,]$ times
the vector of weights. So if we set $r_{jk}=0$ for $j>k$ and collect the weights
as the columns of a matrix $\mt{R}$, then
\[
  \mt{A}=\mt{Q}\mt{R},
  \qquad
  \mt{R}=\begin{bmatrix}r_{11}&r_{12}&\cdots&r_{1n}\\
                        0&r_{22}&\cdots&r_{2n}\\
                        \vdots&&\ddots&\vdots\\
                        0&0&\cdots&r_{nn}\end{bmatrix}.
\]
The zeros below the diagonal are the statement that $\vc{a}_k$ uses only
$\vc{q}_1,\dots,\vc{q}_k$.

%------------------------------------------------------------------------------
\subsection{The theorem}
%------------------------------------------------------------------------------
\begin{thm}[$QR$ decomposition]
\label{thm:qr}
Every $\mt{A}\in\R^{m\times n}$ with independent columns can be factored as
\[
  \boxed{\;\mt{A}=\mt{Q}\mt{R}\;}
\]
where $\mt{Q}\in\R^{m\times n}$ has orthonormal columns ($\mt{Q}\T\mt{Q}=\mt{I}_n$)
and $\mt{R}\in\R^{n\times n}$ is upper triangular with positive diagonal entries
$r_{kk}=\norm{\vc{w}_k}$.
\end{thm}

\noindent With the positive-diagonal requirement the factorization is unique; we
state this without proof. There is also a direct formula for $\mt{R}$. Multiply
$\mt{A}=\mt{Q}\mt{R}$ on the left by $\mt{Q}\T$ and use $\mt{Q}\T\mt{Q}=\mt{I}$:
\[
  \mt{R}=\mt{Q}\T\mt{A},\qquad r_{jk}=\vc{q}_j\T\vc{a}_k .
\]
For $j<k$ this is the shadow length of the algorithm. For $j=k$ it gives
$\vc{q}_k\T\vc{a}_k=\vc{q}_k\T\vc{w}_k=\norm{\vc{w}_k}$. For $j>k$ it is $0$,
because $\vc{a}_k$ lies in the span of $\vc{q}_1,\dots,\vc{q}_k$, all orthogonal to
$\vc{q}_j$.

%------------------------------------------------------------------------------
\subsection{A square example}
%------------------------------------------------------------------------------
\begin{ex}
\label{ex:qrsq}
For the matrix of Example~\ref{ex:gs}, the numbers computed there give
\[
  \mt{Q}=\begin{bmatrix}\frac{1}{\sqrt2}&\frac{1}{\sqrt6}&-\frac{1}{\sqrt3}\\[3pt]
                        \frac{1}{\sqrt2}&-\frac{1}{\sqrt6}&\frac{1}{\sqrt3}\\[3pt]
                        0&\frac{2}{\sqrt6}&\frac{1}{\sqrt3}\end{bmatrix},
  \qquad
  \mt{R}=\begin{bmatrix}\sqrt2&\frac{1}{\sqrt2}&\frac{1}{\sqrt2}\\[3pt]
                        0&\frac{\sqrt6}{2}&\frac{1}{\sqrt6}\\[3pt]
                        0&0&\frac{2}{\sqrt3}\end{bmatrix}.
\]
Check the third column:
$\frac{1}{\sqrt2}\vc{q}_1+\frac{1}{\sqrt6}\vc{q}_2+\frac{2}{\sqrt3}\vc{q}_3
=\frac12[\,1,1,0\,]+\frac16[\,1,-1,2\,]+\frac23[\,-1,1,1\,]=[\,0,1,1\,]=\vc{a}_3$.
\checkmark
\end{ex}

\noindent When $\mt{A}$ is square, $\mt{Q}$ is square with orthonormal columns, so
it is an orthogonal matrix (\S20.6, \S51.3): $\mt{Q}\T\mt{Q}=\mt{Q}\mt{Q}\T=\mt{I}$.
In the language of Part 3, \emph{every invertible matrix is a rigid motion (a
rotation, possibly with a reflection) applied after an upper triangular matrix}.

%------------------------------------------------------------------------------
\subsection{A tall example}
%------------------------------------------------------------------------------
\begin{ex}
\label{ex:qrtall}
The line-fit matrix of \S49.4, $\mt{A}=\begin{bmatrix}1&1\\1&2\\1&3\end{bmatrix}$.
\begin{itemize}[leftmargin=1.4em,itemsep=3pt]
  \item $r_{11}=\norm{[\,1,1,1\,]}=\sqrt3$, $\vc{q}_1=\frac{1}{\sqrt3}[\,1,1,1\,]$.
  \item $r_{12}=\vc{q}_1\T\vc{a}_2=\frac{6}{\sqrt3}=2\sqrt3$, and
        $\vc{w}_2=[\,1,2,3\,]-2\sqrt3\cdot\frac{1}{\sqrt3}[\,1,1,1\,]=[\,-1,0,1\,]$,
        so $r_{22}=\sqrt2$ and $\vc{q}_2=\frac{1}{\sqrt2}[\,-1,0,1\,]$.
\end{itemize}
\[
  \mt{A}=\underbrace{\begin{bmatrix}\frac{1}{\sqrt3}&-\frac{1}{\sqrt2}\\[3pt]
                                     \frac{1}{\sqrt3}&0\\[3pt]
                                     \frac{1}{\sqrt3}&\frac{1}{\sqrt2}\end{bmatrix}}_{\mt{Q},\ 3\times2}
  \underbrace{\begin{bmatrix}\sqrt3&2\sqrt3\\0&\sqrt2\end{bmatrix}}_{\mt{R},\ 2\times2}.
\]
\end{ex}

\noindent Here $\mt{Q}$ is $3\times2$. Its columns are orthonormal, so
$\mt{Q}\T\mt{Q}=\mt{I}_2$, but it is not square and $\mt{Q}\mt{Q}\T\neq\mt{I}_3$
(the warning of \S51.3). In fact
$\mt{Q}\mt{Q}\T=\vc{q}_1\vc{q}_1\T+\vc{q}_2\vc{q}_2\T$ is symmetric and idempotent,
an orthogonal projection, onto $\colsp{\mt{A}}$, as \S57.3 shows. This
$m\times n$ version is sometimes called the \emph{thin} or \emph{reduced} $QR$.
Extending $\vc{q}_1,\dots,\vc{q}_n$ to an orthonormal basis of all of $\R^{m}$
gives a square orthogonal $\mt{Q}$ and the \emph{full} $QR$, with $\mt{R}$ padded
by rows of zeros.

%------------------------------------------------------------------------------
\subsection{The factorization that was promised}
%------------------------------------------------------------------------------
The computational unit ended its Gram--Schmidt part by observing that the
relationship between the input vectors and the output vectors ``is itself a
factorisation of the matrix whose columns are the $\vc{v}_{k}$'', and that naming
it needed matrix machinery not yet available. This is that factorization, and
its name is $QR$.

%==============================================================================
\section{\texorpdfstring{$\colsp{\mt{Q}}=\colsp{\mt{A}}$}{col[Q] = col[A]}}
%==============================================================================

%------------------------------------------------------------------------------
\subsection{The theorem}
%------------------------------------------------------------------------------
\begin{thm}
\label{thm:samecol}
If $\mt{A}=\mt{Q}\mt{R}$ is the $QR$ decomposition of a matrix with independent
columns, then $\colsp{\mt{Q}}=\colsp{\mt{A}}$.
\end{thm}

\begin{proof}
\emph{$\colsp{\mt{A}}\subseteq\colsp{\mt{Q}}$.} Every vector of the column space of
$\mt{A}$ is $\mt{A}\vc{x}=\mt{Q}(\mt{R}\vc{x})$, which is $\mt{Q}$ times a vector.

\emph{$\colsp{\mt{Q}}\subseteq\colsp{\mt{A}}$.} $\mt{R}$ is triangular with nonzero
diagonal entries, so $\det\mt{R}=r_{11}\cdots r_{nn}\neq0$ (\S39.1) and $\mt{R}$ is
invertible (\S51.2). So $\mt{Q}=\mt{A}\mt{R}^{-1}$, and every $\mt{Q}\vc{y}$ equals
$\mt{A}(\mt{R}^{-1}\vc{y})$.
\end{proof}

%------------------------------------------------------------------------------
\subsection{Nested spans}
%------------------------------------------------------------------------------
More is true, one column at a time:
$\spanop\{\vc{q}_1,\dots,\vc{q}_k\}=\spanop\{\vc{a}_1,\dots,\vc{a}_k\}$ for each
$k$. Each $\vc{a}_k$ uses only $\vc{q}_1,\dots,\vc{q}_k$, because $\mt{R}$ is upper
triangular. Each $\vc{q}_k$ uses only $\vc{a}_1,\dots,\vc{a}_k$, because
$\mt{R}^{-1}$ is upper triangular too (\S53.2). This is the property the
process was designed to have, now read off the shape of $\mt{R}$.

%------------------------------------------------------------------------------
\subsection{What it gives}
%------------------------------------------------------------------------------
\begin{itemize}[leftmargin=1.4em,itemsep=4pt]
  \item The columns of $\mt{Q}$ are an \emph{orthonormal basis} of $\colsp{\mt{A}}$.
        Coordinates against an orthonormal basis need no solving, only dot
        products (the computational unit, \S16): the coordinates of
        $\vc{y}\in\colsp{\mt{A}}$ are $\mt{Q}\T\vc{y}$.
  \item The orthogonal projection onto $\colsp{\mt{A}}$ takes its simplest form.
        By \S54.2 it is $\mt{A}(\mt{A}\T\mt{A})^{-1}\mt{A}\T$. Substituting
        $\mt{A}=\mt{Q}\mt{R}$ and $\mt{A}\T\mt{A}=\mt{R}\T\mt{R}$ (\S58.1),
        \[
          \mt{Q}\mt{R}\,(\mt{R}\T\mt{R})^{-1}\mt{R}\T\mt{Q}\T
          =\mt{Q}\mt{R}\,\mt{R}^{-1}(\mt{R}\T)^{-1}\mt{R}\T\mt{Q}\T
          =\mt{Q}\mt{Q}\T .
        \]
        With $n=1$ this is $\hat{\vc{u}}\hat{\vc{u}}\T$, the projection onto a
        line of \S26.2.
  \item If the columns of $\mt{A}$ are \emph{dependent}, some $\vc{w}_k$ is
        $\vc{0}$ and $r_{kk}=0$ (the computational unit, \S19). Then
        $\colsp{\mt{A}}$ has dimension less than $n$ and cannot have $n$
        orthonormal basis vectors, and the theorem does not apply.
\end{itemize}

%==============================================================================
\section{\texorpdfstring{$QR$}{QR} and least squares}
%==============================================================================

%------------------------------------------------------------------------------
\subsection{The normal equations collapse}
%------------------------------------------------------------------------------
With $\mt{A}=\mt{Q}\mt{R}$ and $\mt{Q}\T\mt{Q}=\mt{I}$,
\[
  \mt{A}\T\mt{A}=\mt{R}\T\mt{Q}\T\mt{Q}\mt{R}=\mt{R}\T\mt{R},
  \qquad
  \mt{A}\T\vc{b}=\mt{R}\T\mt{Q}\T\vc{b}.
\]
So the normal equations $\mt{A}\T\mt{A}\xh=\mt{A}\T\vc{b}$ become
$\mt{R}\T\mt{R}\xh=\mt{R}\T\mt{Q}\T\vc{b}$. The matrix $\mt{R}\T$ is invertible
(its determinant is $\det\mt{R}\neq0$), so multiply both sides by
$(\mt{R}\T)^{-1}$:

\begin{thm}
\label{thm:qrls}
If $\mt{A}=\mt{Q}\mt{R}$, the least-squares solution of $\mt{A}\vc{x}=\vc{b}$ is the
solution of the triangular system
\[
  \boxed{\;\mt{R}\,\xh=\mt{Q}\T\vc{b},\;}
\]
found by one back substitution (\S41).
\end{thm}

\noindent The pseudo-inverse simplifies the same way:
$\mt{A}^{+}=(\mt{A}\T\mt{A})^{-1}\mt{A}\T=(\mt{R}\T\mt{R})^{-1}\mt{R}\T\mt{Q}\T
=\mt{R}^{-1}(\mt{R}\T)^{-1}\mt{R}\T\mt{Q}\T=\mt{R}^{-1}\mt{Q}\T$.

%------------------------------------------------------------------------------
\subsection{The line fit, a third time}
%------------------------------------------------------------------------------
\begin{ex}
With $\mt{Q}$ and $\mt{R}$ from Example~\ref{ex:qrtall} and $\vc{b}=[\,1,2,2\,]$:
\[
  \mt{Q}\T\vc{b}=\begin{bmatrix}\frac{1+2+2}{\sqrt3}\\[3pt]\frac{-1+0+2}{\sqrt2}\end{bmatrix}
  =\begin{bmatrix}\frac{5}{\sqrt3}\\[3pt]\frac{1}{\sqrt2}\end{bmatrix},
  \qquad
  \begin{bmatrix}\sqrt3&2\sqrt3\\0&\sqrt2\end{bmatrix}
  \begin{bmatrix}c_0\\c_1\end{bmatrix}
  =\begin{bmatrix}\frac{5}{\sqrt3}\\[3pt]\frac{1}{\sqrt2}\end{bmatrix}.
\]
Back substitution: $c_1=\frac{1}{\sqrt2}\big/\sqrt2=\frac12$, then
$c_0=\bigl(\frac{5}{\sqrt3}-2\sqrt3\cdot\frac12\bigr)\big/\sqrt3=\frac53-1=\frac23$.
The line $y=\frac23+\frac12t$ of \S49.4 and \S53.4, found a third way.
\end{ex}

%------------------------------------------------------------------------------
\subsection{Three routes to one answer}
%------------------------------------------------------------------------------
\begin{center}
\renewcommand{\arraystretch}{1.4}
\begin{tabular}{@{}l l l@{}}
\toprule
\textbf{Route} & \textbf{Solve} & \textbf{Where} \\
\midrule
normal equations, by elimination & $\mt{A}\T\mt{A}\xh=\mt{A}\T\vc{b}$ & \S49 \\
normal equations, by the inverse & $\xh=(\mt{A}\T\mt{A})^{-1}\mt{A}\T\vc{b}=\mt{A}^{+}\vc{b}$ & \S53.4, \S54 \\
$QR$ & $\mt{R}\xh=\mt{Q}\T\vc{b}$, back substitution & \S58.1 \\
\bottomrule
\end{tabular}
\end{center}

\noindent In exact arithmetic all three give the same $\xh$. On a computer they do
not behave the same, and that is the reason for the third.

%------------------------------------------------------------------------------
\subsection{Why \texorpdfstring{$QR$}{QR} in practice}
%------------------------------------------------------------------------------
Forming $\mt{A}\T\mt{A}$ is the weak step of the normal equations. When the
columns of $\mt{A}$ are \emph{nearly} dependent, as often happens with real data,
$\mt{A}$ is already sensitive to small errors, and $\mt{A}\T\mt{A}$ is much more
so: roughly, its sensitivity is the square of $\mt{A}$'s. A problem that could be
solved to eight correct digits may come back with four. The $QR$ route never forms
$\mt{A}\T\mt{A}$. It works with orthonormal columns, and multiplying by an
orthogonal matrix does not stretch anything (\S20), so it does not stretch errors
either.

Two refinements matter for large problems. The classical process of \S55 slowly
loses orthogonality in floating point; the \emph{modified} Gram--Schmidt process
rearranges the same arithmetic and holds up far better (the computational unit,
\S20.3). And professional software usually computes $QR$ by a different method
altogether, with a sequence of reflections $\mt{I}-2\mt{P}_{\vc{v}}$, the
Householder reflectors of \S28.5--\S28.7. Each one zeros out a
column below the diagonal while keeping lengths exactly.

%------------------------------------------------------------------------------
\subsection{Practice}
%------------------------------------------------------------------------------
\begin{enumerate}[label=\textbf{P\arabic*.},leftmargin=2.6em,itemsep=5pt]
  \item Find the $QR$ decomposition of $\begin{bmatrix}3&1\\4&2\end{bmatrix}$, and
        check that $\det\mt{A}=\det\mt{Q}\,\det\mt{R}$.
  \item Find the (thin) $QR$ decomposition of
        $\mt{A}=\begin{bmatrix}1&0\\1&1\\0&1\end{bmatrix}$.
  \item Use P2 to find the least-squares solution of $\mt{A}\vc{x}=[\,1,0,1\,]$, and
        check it against the normal equations.
  \item Show that for a square $\mt{A}=\mt{Q}\mt{R}$,
        $|\det\mt{A}|=r_{11}r_{22}\cdots r_{nn}=\norm{\vc{w}_1}\cdots\norm{\vc{w}_n}$,
        and interpret this for $n=2$ as base times height.
  \item Explain why $\mt{Q}\mt{Q}\T$ in Example~\ref{ex:qrtall} cannot be
        $\mt{I}_3$.
\end{enumerate}

\subsubsection*{Answers}

\noindent\textbf{P1.} $r_{11}=5$, $\vc{q}_1=[\,\frac35,\frac45\,]$;
$r_{12}=\frac{3+8}{5}=\frac{11}{5}$;
$\vc{w}_2=[\,1,2\,]-\frac{11}{5}[\,\frac35,\frac45\,]=[\,-\frac{8}{25},\frac{6}{25}\,]$,
$r_{22}=\frac25$, $\vc{q}_2=[\,-\frac45,\frac35\,]$. So
$\mt{Q}=\begin{bmatrix}\frac35&-\frac45\\[2pt]\frac45&\frac35\end{bmatrix}$ (a rotation,
$\det=1$) and $\mt{R}=\begin{bmatrix}5&\frac{11}{5}\\[2pt]0&\frac25\end{bmatrix}$
($\det=2$); and $\det\mt{A}=6-4=2$. \checkmark

\smallskip
\noindent\textbf{P2.} $r_{11}=\sqrt2$, $\vc{q}_1=\frac{1}{\sqrt2}[\,1,1,0\,]$;
$r_{12}=\frac{1}{\sqrt2}$; $\vc{w}_2=[\,0,1,1\,]-\frac12[\,1,1,0\,]=[\,-\frac12,\frac12,1\,]$,
$r_{22}=\sqrt{\frac32}$, $\vc{q}_2=\frac{1}{\sqrt6}[\,-1,1,2\,]$.
$\mt{R}=\begin{bmatrix}\sqrt2&\frac{1}{\sqrt2}\\[2pt]0&\sqrt{\frac32}\end{bmatrix}$.

\smallskip
\noindent\textbf{P3.} $\mt{Q}\T\vc{b}=[\,\frac{1}{\sqrt2},\frac{1}{\sqrt6}\,]$.
Back substitution: $x_2=\frac{1}{\sqrt6}\big/\sqrt{\frac32}=\frac13$, and
$x_1=\bigl(\frac{1}{\sqrt2}-\frac{1}{\sqrt2}\cdot\frac13\bigr)\big/\sqrt2=\frac13$.
Normal equations: $\mt{A}\T\mt{A}=\begin{bmatrix}2&1\\1&2\end{bmatrix}$,
$\mt{A}\T\vc{b}=[\,1,1\,]$, solved by $[\,\frac13,\frac13\,]$. \checkmark

\smallskip
\noindent\textbf{P4.} $\det\mt{A}=\det\mt{Q}\det\mt{R}$ with $\det\mt{Q}=\pm1$
(\S39.3) and $\det\mt{R}=r_{11}\cdots r_{nn}$ (\S39.1). For $n=2$:
$r_{11}=\norm{\vc{a}_1}$ is the base, and $r_{22}=\norm{\vc{w}_2}$ is the length of
the part of $\vc{a}_2$ perpendicular to $\vc{a}_1$, the height. So this is
\S23's base times height, computed by Gram--Schmidt.

\smallskip
\noindent\textbf{P5.} $\mt{Q}\mt{Q}\T$ has rank at most $2$ (its column space is
inside $\colsp{\mt{Q}}$, a plane), while $\mt{I}_3$ has rank $3$.

\lookahead This ends the matrix unit. Its last theorems were about choosing a good
basis: an orthonormal basis of the column space makes least squares easy. The next
unit asks for a basis that makes a square matrix itself easy. The ideal basis
vectors are the ones the matrix only stretches, its \emph{eigenvectors}, first met
as the invariant lines of \S30.4.

\end{document}
